\documentclass[12pt]{article}
\usepackage{enumitem}
\usepackage{geometry}
\usepackage{amsmath, amssymb}
\geometry{margin=1in}
\begin{document}
\begin{center}
Chemistry - K-12 Level Assessment 1 \
Total Marks: 40 \
Answer all questions.
\end{center}

\section*{Multiple Choice (20 marks)}
\begin{enumerate}[label=\arabic*.]
\item A 2.35-gram sample was dissolved in water, and the chloride ions were precipitated by adding silver nitrate (Ag+ + Cl- $\rightarrow$ AgCl). If 0.435 g of precipitate was...
\begin{enumerate}[label=(\alph*)]
    \item 10.80\%
    \item 4.60\%
    \item 43.50\%
    \item 18\%
\end{enumerate}
\item A 2.0 L flask holds 0.40 g of helium gas. If the helium is evacuated into a larger container while the temperature is held constant, what will the effect on...
\begin{enumerate}[label=(\alph*)]
    \item It will remain constant as the number of helium molecules does not change.
    \item It will decrease as the gas will be more ordered in the larger flask.
    \item It will decrease because the molecules will collide with the sides of the larger flask less often than they did in the smaller flask.
    \item It will increase as the gas molecules will be more dispersed in the larger flask.
\end{enumerate}
\item The electrolysis of molten magnesium bromide is expected to produce
\begin{enumerate}[label=(\alph*)]
    \item magnesium at the anode and bromine at the cathode
    \item magnesium at the cathode and bromine at the anode
    \item magnesium at the cathode and oxygen at the anode
    \item bromine at the anode and hydrogen at the cathode
\end{enumerate}
\item What is the general relationship between temperature and entropy for diatomic gases?
\begin{enumerate}[label=(\alph*)]
    \item They are completely independent of each other; temperature has no effect on entropy.
    \item There is a direct relationship, because at higher temperatures there is an increase in energy dispersal.
    \item There is an inverse relationship, because at higher temperatures substances are more likely to be in a gaseous state.
    \item It depends on the specific gas and the strength of the intermolecular forces between individual molecules.
\end{enumerate}
\item NH4+(aq) + NO2- (aq) $\rightarrow$ N2(g) + 2H2O(l) Increasing the temperature of the above reaction will increase the rate of reaction. Which of the following is NOT a...
\begin{enumerate}[label=(\alph*)]
    \item The reactants will be more likely to overcome the activation energy.
    \item The number of collisions between reactant molecules will increase.
    \item A greater distribution of reactant molecules will have high velocities.
    \item Alternate reaction pathways become available at higher temperatures.
\end{enumerate}
\item Three half-lives after an isotope is prepared,
\begin{enumerate}[label=(\alph*)]
    \item 25\% of the isotope is left
    \item 25\% of the isotope decayed
    \item 12.5\% of the isotope is left
    \item 12.5\% of the isotope decayed
\end{enumerate}
\item Which of the following processes is an irreversible reaction?
\begin{enumerate}[label=(\alph*)]
    \item CH4(g) + O2(g) $\rightarrow$ CO2(g) + H2O(l)
    \item HCN(aq) + H2O(l) $\rightarrow$ CN-(aq) + H3O+(aq)
    \item Al(NO3)3(s) $\rightarrow$ Al3+(aq) + 3NO3-(aq)
    \item 2Ag+(aq) + Ti(s) $\rightarrow$ 2Ag(s) + Ti2+(aq)
\end{enumerate}
\item When potassium perchlorate, KClO4, dissolves in water, the temperature of the resultant solution is lower than the initial temperature of the components....
\begin{enumerate}[label=(\alph*)]
    \item This is a spontaneous process because it is exothermic.
    \item This is a spontaneous process because of an entropy increase.
    \item This is a spontaneous process because of an entropy decrease.
    \item This is a spontaneous process because it is exothermic.
\end{enumerate}
\item When an ideal gas is allowed to expand isothermally, which one of the following is true?
\begin{enumerate}[label=(\alph*)]
    \item q = 0
    \item w = 0
    \item E = 0
    \item q = -w
\end{enumerate}
\item The rate of a chemical reaction is determined by
\begin{enumerate}[label=(\alph*)]
    \item the equilibrium constant
    \item the rate-determining or slow step of the mechanism
    \item the reaction vessel pressure
    \item the intermediates formed in the first step
\end{enumerate}
\end{enumerate}

\section*{True / False (20 marks)}
\textit{Answer True or False and justify briefly.}
\begin{enumerate}[label=\arabic*.]
\item True or False: The correct answer to 'Sulfurous acid is a weak acid, while sulfuric acid is a much stronger acid because' is 'sulfuric acid has more oxygen atoms in its formula'.
\item True or False: The correct answer to 'Chemical reactions can be classified as either heterogeneous or homogeneous. Which of the following equations below is best classified as...' is 'C(s) + H2O(g) $\rightarrow$ H2(g) + CO(g)'.
\item True or False: The correct answer to 'SO2Cl2 $\rightarrow$ SO2(g) + Cl2(g) At 600 K, SO2Cl2 will decompose to form sulfur dioxide and chlorine gas via the above equation. If the reaction...' is 'Increasing the temperature at which the reaction occurs'.
\item True or False: The correct answer to 'S(s) + O2(g) $\rightarrow$ SO2(g)' is 'Q must be close to 1.0 since there is one mole of gas on each side of the equation'.
\item True or False: The correct answer to 'What is the molarity of a sodium hydroxide solution that requires 42.6 mL of 0.108 M HCl to neutralize 40.0 mL of the base?' is '0.115 M'.
\item True or False: The correct answer to 'Each resonance form of the nitrate ion, NO3-, has how many sigma and how many pi bonds?' is '1 sigma and 1 pi'.
\item True or False: The correct answer to 'S(s) + O2(g) $\rightarrow$ SO2(g)' is 'one mole of sulfur atoms reacts with one mole of oxygen molecules to yield one mole of sulfur dioxide molecules'.
\item True or False: The correct answer to 'Approximately how many half-lives will it take to reduce 10.0 kg of a radioactive substance to 1.0 microgram of that substance?' is '13'.
\item True or False: The correct answer to 'When the dichromate ion is reacted, one of its most common products is Cr3+. What is the oxidation state (oxidation number) of chromium...' is '6+, reduction'.
\item True or False: The correct answer to 'How many electrons, neutrons, and protons are in an atom of Cr?' is '27 electrons, 27 protons, 24 neutrons'.
\end{enumerate}

\vfill
\noindent\textit{End of Paper}
\end{document}